\documentclass{beamer}

\usepackage{beamerthemeSingapore}
\usepackage[brazil]{babel}
\usepackage[T1]{fontenc}
\usepackage[utf8]{inputenc}
\usepackage{graphicx}
\usepackage{url}
\usepackage{xcolor}
\usepackage{booktabs}
\usepackage{ragged2e}
\usepackage{etoolbox}
\usepackage{minted}
\usepackage{array}

\newcommand{\tcb}[1]{\textcolor{blue}{#1}}
\newcommand{\tcr}[1]{\textcolor{red}{#1}}
\newcommand{\tcg}[1]{\textcolor{green!50!black}{#1}}

\setminted{
  fontsize=\footnotesize,
  breaklines=true,
  breakanywhere=true,
  frame=lines,
  framesep=2mm,
  baselinestretch=1.05,
  tabsize=2
}

\apptocmd{\frame}{}{\justifying}{}

\setbeamertemplate{navigation symbols}{}

% slides de prática e de resposta: fundo e título em cores próprias
\newcommand{\estilopratica}{%
  \setbeamercolor{frametitle}{fg=orange!85!black}%
  \setbeamercolor{background canvas}{bg=orange!8}}
\newcommand{\estiloresposta}{%
  \setbeamercolor{frametitle}{fg=green!45!black}%
  \setbeamercolor{background canvas}{bg=green!6}}

\title{11 - Back-end e PHP}
\author{Prof. Alex Kutzke}
\date{Outubro de 2026}

\begin{document}

% ==========================================
\begin{frame}
	\begin{center}
		\large{Setor de Educação Profissional e Tecnológica}\\
		Tecnologia em Análise e Desenvolvimento de Sistemas\\
		\vspace{1cm}
		\small{DS122 - Desenvolvimento de Aplicações Web 1}
	\end{center}
	\titlepage
\end{frame}

% ==========================================
\begin{frame}{Objetivos da aula}
	Ao final desta aula você deve ser capaz de:
	\begin{enumerate}
		\item Explicar o que o interpretador do PHP faz com um arquivo
		      \texttt{.php}, e o que chega ao navegador;
		\item Iniciar o servidor embutido do PHP e abrir uma página por ele;
		\item Declarar variáveis e usar \texttt{==} e \texttt{===};
		\item Montar textos com concatenação, interpolação e funções de string;
		\item Criar arrays indexados e associativos, e percorrê-los com
		      \texttt{foreach};
		\item Usar \texttt{if}, \texttt{elseif} e \texttt{else} para decidir que
		      HTML é produzido;
		\item Ler um arquivo JSON com \texttt{json\_decode} e montar o catálogo a
		      partir dele.
	\end{enumerate}
\end{frame}

% ==========================================
\begin{frame}{Roteiro}
	\tableofcontents
\end{frame}

% ==========================================
\section{Back-end}
% ==========================================

\begin{frame}[fragile]{Uma página gerada no servidor}
	\texttt{ola.php}:

	\begin{minted}{html+php}
<!DOCTYPE html>
<html lang="pt-BR">
<head>
  <meta charset="UTF-8">
  <title>Olá, PHP</title>
</head>
<body>
  <h1>Olá, PHP</h1>
  <p>Número sorteado no servidor: <?= rand(1, 6) ?></p>
</body>
</html>
	\end{minted}

	No terminal, na pasta do arquivo:

	\begin{minted}{bash}
php -d display_errors=1 -S localhost:8000
	\end{minted}

	\quad e no navegador: \texttt{http://localhost:8000/ola.php}
\end{frame}

\begin{frame}[fragile]{O que chegou ao navegador}
	Cada recarga sorteia outro número.

	\vspace{0.2cm}
	Código-fonte da página, com \texttt{Ctrl+U}:

	\begin{minted}{html}
<body>
  <h1>Olá, PHP</h1>
  <p>Número sorteado no servidor: 4</p>
</body>
	\end{minted}

	\vspace{0.2cm}
	Nenhuma linha de PHP. O servidor passou o arquivo ao interpretador do PHP,
	que executou o trecho entre \texttt{<?=} e \texttt{?>} e devolveu o HTML
	pronto. A resposta HTTP levou só o HTML.
\end{frame}

\begin{frame}{Back-end}
	\textbf{Back-end} (\emph{server-side}): a parte da aplicação Web executada no
	servidor. Processa as regras de negócio, produz o HTML com conteúdo dinâmico e
	guarda os dados.

	\vspace{0.3cm}
	Para gerar páginas no servidor, três peças:
	\begin{itemize}
		\item um \textbf{servidor web}, que recebe as requisições HTTP e envia as
		      respostas;
		\item um \textbf{interpretador} (ou programa), que executa a lógica e
		      produz o HTML;
		\item um \textbf{banco de dados}, que guarda os dados da aplicação.
	\end{itemize}

	\vspace{0.3cm}
	O banco de dados entra na disciplina na aula de PHP e MySQL.
\end{frame}

\begin{frame}{Servidor web e diretório raiz}
	O servidor web escuta requisições, em geral na porta 80 (HTTP) e 443 (HTTPS).
	Exemplos: Apache, Nginx, Lighttpd.

	\vspace{0.3cm}
	Em páginas estáticas, ele traduz a URL para um arquivo dentro do
	\textbf{diretório raiz} (\emph{document root}):

	\vspace{0.2cm}
	\begin{center}
		\small
		\texttt{http://example.com/dir1/arquivo.html}\\[0.1cm]
		\texttt{/var/www/html/dir1/arquivo.html}
	\end{center}

	\vspace{0.2cm}
	O servidor nega acesso a arquivos fora do diretório raiz.
\end{frame}

\begin{frame}{Páginas dinâmicas}
	Num sistema de notas, cada aluno vê as suas. Um arquivo HTML por aluno é
	inviável: o HTML precisa ser gerado a cada requisição.

	\vspace{0.3cm}
	\begin{enumerate}
		\item o navegador pede \texttt{/catalogo.php};
		\item o servidor web vê o \texttt{.php} e passa o arquivo ao interpretador;
		\item o interpretador executa o código e produz o HTML;
		\item o servidor envia esse HTML na resposta.
	\end{enumerate}

	\vspace{0.3cm}
	Quem executa o código é o interpretador, não o servidor web. No Apache, por
	exemplo, a ligação é feita pelo módulo \texttt{mod\_php} ou pelo serviço
	PHP-FPM.
\end{frame}

\begin{frame}{Renderização no servidor, APIs e SPAs}
	Nesta parte da disciplina: \textbf{renderização no servidor} (SSR). O
	back-end monta o HTML completo de cada página.

	\vspace{0.3cm}
	Outro modelo, já visto pela metade com \texttt{fetch} e
	\texttt{produtos.json}:
	\begin{itemize}
		\item \textbf{API}: o back-end devolve apenas os dados, em geral em JSON,
		      gerados a cada requisição;
		\item \textbf{SPA}: o front-end em JavaScript pede os dados com
		      \texttt{fetch} e monta as telas pelo DOM, sem recarregar a página.
	\end{itemize}

	\vspace{0.3cm}
	O segundo modelo é assunto de disciplinas seguintes do curso.
\end{frame}

% ==========================================
\section{Servidor}
% ==========================================

\begin{frame}[fragile]{O servidor embutido do PHP}
	\begin{minted}{bash}
php -d display_errors=1 -S localhost:8000
	\end{minted}

	\begin{itemize}
		\item serve os arquivos da pasta onde o comando foi executado;
		\item \texttt{.php} passa pelo interpretador; HTML, CSS, imagens e JSON
		      saem como estão;
		\item endereço sem arquivo: procura \texttt{index.php}, depois
		      \texttt{index.html};
		\item \texttt{Ctrl+C} encerra.
	\end{itemize}

	\vspace{0.2cm}
	Feito para desenvolvimento. O manual do PHP avisa para não usá-lo em rede
	pública.

	\vspace{0.2cm}
	Instalação: \texttt{php-cli} pelo \texttt{apt} ou \texttt{dnf}; no Windows,
	o PHP do XAMPP. Ver o material de Programação Back-end.
\end{frame}

\begin{frame}[fragile]{Erros: página e terminal}
	\texttt{-d display\_errors=1}: a mensagem de erro aparece na página. Sem a
	opção, conforme a instalação, a página chega em branco.

	\vspace{0.2cm}
	O terminal do servidor mostra uma linha por requisição, com status e erro:

	\begin{minted}[fontsize=\scriptsize]{text}
[Mon Oct 26 12:35:47 2026] [::1]:33222 [200]: GET /ola.php
[Mon Oct 26 12:35:54 2026] [::1]:33250 [404]: GET /nada.php
                                   - No such file or directory
	\end{minted}

	\vspace{0.2cm}
	Live Server, \texttt{python3 -m http.server} e \texttt{file://} entregam o
	\texttt{.php} sem executar: o navegador mostra o código ou baixa o arquivo.
\end{frame}

{\estilopratica
\begin{frame}[fragile]{Prática: o ambiente}
	\begin{enumerate}
		\item confira a instalação com \texttt{php -v};
		\item crie a pasta \texttt{aula-php/}, fora do repositório do trabalho;
		\item crie nela o \texttt{ola.php};
		\item inicie o servidor dentro da pasta e abra
		      \texttt{http://localhost:8000/ola.php};
		\item recarregue algumas vezes e abra o código-fonte com
		      \texttt{Ctrl+U}.
	\end{enumerate}

	\vspace{0.3cm}
	Os exercícios desta aula não têm entrega.

	\vspace{0.3cm}
	Tempo: 10 minutos.
\end{frame}
}

{\estiloresposta
\begin{frame}[fragile]{Resposta: o ambiente}
	\begin{minted}{text}
$ php -v
PHP 8.3.x (cli) ...
	\end{minted}

	\begin{itemize}
		\item o número muda a cada recarga;
		\item o código-fonte mostra só o número, sem \texttt{<?=};
		\item no terminal, uma linha \texttt{[200]: GET /ola.php} por recarga.
	\end{itemize}

	\vspace{0.2cm}
	\texttt{php: command not found}: o PHP não está instalado, ou não está no
	\texttt{PATH}.
\end{frame}
}

% ==========================================
\section{Sintaxe}
% ==========================================

\begin{frame}[fragile]{Tags PHP}
	O interpretador executa o que está entre as tags. O resto vai para a saída
	sem alteração.

	\vspace{0.2cm}
	\begin{itemize}
		\item \texttt{<?php ... ?>}: bloco de código;
		\item \texttt{<?= ... ?>}: imprime um valor; equivale a
		      \texttt{<?php echo ... ?>}.
	\end{itemize}

	\begin{minted}{html+php}
<p>Isto vai direto ao navegador.</p>
<?php echo 'Isto é processado no servidor.'; ?>
<p>Isto também vai direto.</p>
	\end{minted}

	O navegador recebe:

	\begin{minted}{html}
<p>Isto vai direto ao navegador.</p>
Isto é processado no servidor.<p>Isto também vai direto.</p>
	\end{minted}

	A quebra de linha logo depois de \texttt{?>} não vai para a saída.
\end{frame}

\begin{frame}[fragile]{Ponto e vírgula e a tag de fechamento}
	Cada instrução termina com \texttt{;}.

	\vspace{0.2cm}
	Arquivo só com PHP, sem HTML: \textbf{sem} \texttt{?>} no fim. Um espaço
	depois dele iria para a saída e causaria o erro \emph{headers already sent},
	visto nas aulas de formulários e de sessões.

	\begin{minted}{php}
<?php
// Arquivo com apenas código PHP: termina sem a tag de fechamento.

echo 'Isto é um teste';
$variavel = 10;
	\end{minted}

	\vspace{0.2cm}
	Um \texttt{?>} dentro de um comentário \texttt{//} também encerra o modo
	PHP. O que vem depois dele vira texto na página.
\end{frame}

\begin{frame}[fragile]{Variáveis e tipos}
	Toda variável começa com \texttt{\$}. Sem \texttt{let} ou \texttt{const}: ela
	existe a partir da primeira atribuição. \texttt{\$preco} e \texttt{\$Preco}
	são variáveis diferentes.

	\begin{minted}{php}
<?php
$nome = "Café em grão";   // string
$estoque = 12;            // int
$preco = 32.5;            // float
$disponivel = true;       // bool
$cupom = null;            // null

var_dump($preco);   // float(32.5)
var_dump($nome);    // string(14) "Café em grão"
	\end{minted}

	\texttt{14} é o tamanho em bytes: \texttt{é} e \texttt{ã} ocupam dois bytes
	em UTF-8.
\end{frame}

\begin{frame}[fragile]{Operadores e comparação}
	Aritméticos e de comparação iguais aos do JavaScript: \texttt{+ - * / \%},
	\texttt{< > <= >=}, \texttt{\&\& || !}.

	\begin{minted}{php}
var_dump(1 == "1");    // bool(true): converte o tipo
var_dump(1 === "1");   // bool(false): compara valor e tipo
	\end{minted}

	\vspace{0.2cm}
	Como no JavaScript, prefira \texttt{===} e \texttt{!==}.
\end{frame}

\begin{frame}[fragile]{Concatenação com ponto}
	O \texttt{+} é sempre soma. A concatenação é com \texttt{.}:

	\begin{minted}{php}
var_dump("1" + "1");   // int(2)
var_dump("1" . "1");   // string(2) "11"
	\end{minted}

	\texttt{"Preço: " + \$preco} interrompe o script:

	\begin{minted}[fontsize=\scriptsize]{text}
Fatal error: Uncaught TypeError: Unsupported operand types: string + float
	\end{minted}

	\texttt{.=} acrescenta ao fim:

	\begin{minted}{php}
$html = "<ul>";
$html .= "<li>Café em grão</li>";
$html .= "</ul>";
	\end{minted}
\end{frame}

\begin{frame}[fragile]{Aspas simples e aspas duplas}
	Aspas duplas substituem a variável pelo valor (\textbf{interpolação}), como a
	\emph{template string} do JavaScript. Aspas simples não:

	\begin{minted}{php}
$nome = "Café em grão";
echo "Produto: $nome";      // Produto: Café em grão
echo 'Produto: $nome';      // Produto: $nome
echo "Produto: {$nome}!";   // Produto: Café em grão!
	\end{minted}

	\vspace{0.2cm}
	As chaves delimitam a variável. Use-as quando a variável é seguida de letra
	ou é um elemento de array.
\end{frame}

\begin{frame}{Funções de string}
	\footnotesize
	\begin{tabular}{>{\raggedright\arraybackslash}p{0.62\textwidth}>{\raggedright\arraybackslash}p{0.28\textwidth}}
		\toprule
		Chamada & Resultado \\
		\midrule
		\texttt{strlen("Chá")} & \texttt{4} (bytes) \\ \midrule
		\texttt{strtoupper("Bebidas")} & \texttt{"BEBIDAS"} \\ \midrule
		\texttt{strtolower("Bebidas")} & \texttt{"bebidas"} \\ \midrule
		\texttt{trim("\ \ Café\ \ ")} & \texttt{"Café"} \\ \midrule
		\texttt{str\_contains("Café em grão", "grão")} & \texttt{true} \\
		\midrule
		\texttt{str\_replace(".", ",", "32.5")} & \texttt{"32,5"} \\ \midrule
		\texttt{number\_format(1234.5, 2, ",", ".")} & \texttt{"1.234,50"} \\
		\bottomrule
	\end{tabular}

	\vspace{0.3cm}
	\small
	\texttt{number\_format}: número, casas decimais, separador decimal,
	separador de milhar. Arredonda: \texttt{27.625} vira \texttt{"27,63"}.

	\vspace{0.1cm}
	\texttt{strtoupper("Café")} dá \texttt{"CAFé"}: só letras sem acento.
\end{frame}

{\estilopratica
\begin{frame}[fragile]{Prática: o preço com desconto}
	Crie o \texttt{preco.php}, com no topo:

	\begin{minted}{php}
<?php
$nome = "Café em grão";
$categoria = "Bebidas";
$preco = 32.5;
$desconto = 0.15;
?>
	\end{minted}

	\begin{enumerate}
		\item calcule o preço com desconto em \texttt{\$precoFinal};
		\item escreva um \texttt{<h1>} com o nome e a categoria em maiúsculas;
		\item escreva um \texttt{<p>} com os dois preços, com
		      \texttt{number\_format};
		\item troque um \texttt{.} por \texttt{+} e leia o erro; desfaça.
	\end{enumerate}

	Confira: \texttt{Café em grão (BEBIDAS)} e
	\texttt{De R\$ 32,50 por R\$ 27,63}.

	\vspace{0.1cm}
	Tempo: 10 minutos.
\end{frame}
}

{\estiloresposta
\begin{frame}[fragile]{Resposta: o preço com desconto}
	\begin{minted}[fontsize=\scriptsize]{html+php}
<?php
$nome = "Café em grão";
$categoria = "Bebidas";
$preco = 32.5;
$desconto = 0.15;
$precoFinal = $preco * (1 - $desconto);
?>
<!DOCTYPE html>
<html lang="pt-BR">
<head>
  <meta charset="UTF-8">
  <title>Preço</title>
</head>
<body>
  <h1><?= $nome . " (" . strtoupper($categoria) . ")" ?></h1>
  <p>De R$ <?= number_format($preco, 2, ",", ".") ?>
     por R$ <?= number_format($precoFinal, 2, ",", ".") ?></p>
</body>
</html>
	\end{minted}
\end{frame}
}

% ==========================================
\section{Arrays}
% ==========================================

\begin{frame}[fragile]{Arrays indexados}
	Valores em sequência, a partir do índice \texttt{0}:

	\begin{minted}{php}
$categorias = ["Bebidas", "Doces", "Grãos"];
echo $categorias[0];        // Bebidas
echo count($categorias);    // 3

$categorias[] = "Temperos"; // acrescenta ao fim
echo count($categorias);    // 4
	\end{minted}

	\vspace{0.2cm}
	\texttt{\$categorias[] = ...} faz o papel do \texttt{push} do JavaScript.
\end{frame}

\begin{frame}[fragile]{Arrays associativos}
	Cada valor tem uma \textbf{chave} de texto. \texttt{=>} liga a chave ao
	valor:

	\begin{minted}{php}
$produto = [
  "id" => 1,
  "nome" => "Café em grão",
  "categoria" => "Bebidas",
  "preco" => 32.5,
  "desconto" => 0.15,
];

echo $produto["nome"];                // Café em grão
echo "Preço: {$produto["preco"]}";    // Preço: 32.5
	\end{minted}

	O papel do objeto no JavaScript, com acesso sempre por colchetes:
	\texttt{\$produto["nome"]}, e não \texttt{produto.nome}.
\end{frame}

\begin{frame}[fragile]{Array de arrays}
	A forma do \texttt{produtos.js} e do \texttt{produtos.json}:

	\begin{minted}{php}
$produtos = [
  ["nome" => "Café em grão", "preco" => 32.5],
  ["nome" => "Chá mate", "preco" => 12.9],
];

echo $produtos[1]["nome"];   // Chá mate
print_r($produtos[0]);
	\end{minted}

	\begin{minted}[fontsize=\scriptsize]{text}
Array
(
    [nome] => Café em grão
    [preco] => 32.5
)
	\end{minted}

	\texttt{echo \$produtos} imprime só \texttt{Array}, com o aviso
	\texttt{Array to string conversion}.
\end{frame}

% ==========================================
\section{Controle}
% ==========================================

\begin{frame}[fragile]{\texttt{if}, \texttt{elseif} e \texttt{else}}
	\begin{minted}{php}
if ($produto["desconto"] > 0.2) {
  echo "Promoção";
} elseif ($produto["desconto"] > 0) {
  echo "Com desconto";
} else {
  echo "Preço normal";
}
	\end{minted}

	\vspace{0.2cm}
	\texttt{elseif} é uma palavra só.
\end{frame}

\begin{frame}[fragile]{\texttt{foreach}}
	Percorre um array, um elemento por vez:

	\begin{minted}{php}
foreach ($produtos as $produto) {
  echo "<li>{$produto["nome"]}</li>";
}
	\end{minted}

	Com a chave:

	\begin{minted}{php}
foreach ($produto as $chave => $valor) {
  echo "$chave: $valor\n";
}
	\end{minted}

	\begin{minted}[fontsize=\scriptsize]{text}
id: 1
nome: Café em grão
categoria: Bebidas
...
	\end{minted}
\end{frame}

\begin{frame}[fragile]{\texttt{for} e \texttt{while}}
	A forma do JavaScript, com \texttt{\$} nas variáveis:

	\begin{minted}{php}
for ($i = 1; $i <= 3; $i++) {
  echo "<p>Parcela $i</p>";
}

$estoque = 3;
while ($estoque > 0) {
  echo "<p>Restam $estoque</p>";
  $estoque--;
}
	\end{minted}
\end{frame}

\begin{frame}[fragile]{Sintaxe alternativa dentro do HTML}
	\texttt{:} no lugar de \texttt{\{}, e \texttt{endif;},
	\texttt{endforeach;} no lugar de \texttt{\}}:

	\begin{minted}{html+php}
<ul>
  <?php foreach ($produtos as $produto): ?>
    <li>
      <?= $produto["nome"] ?>
      <?php if ($produto["preco"] < 20): ?>
        <strong>até R$ 20</strong>
      <?php endif; ?>
    </li>
  <?php endforeach; ?>
</ul>
	\end{minted}

	O \texttt{<li>} se repete por produto; o \texttt{<strong>} só sai quando a
	condição é verdadeira. O HTML fica fora das strings.
\end{frame}

{\estilopratica
\begin{frame}[fragile]{Prática: o catálogo com \texttt{foreach}}
	Crie o \texttt{catalogo.php} com o array de três produtos da seção 14 do
	material (\texttt{id}, \texttt{nome}, \texttt{descricao},
	\texttt{categoria}, \texttt{preco}, \texttt{desconto}).

	\begin{enumerate}
		\item escreva um \texttt{<p>} com a quantidade de produtos, com
		      \texttt{count};
		\item escreva um \texttt{<article class="cartao">} por produto, com
		      \texttt{foreach} na sintaxe alternativa: \texttt{<h2>} com o nome,
		      \texttt{<p>} com a descrição, \texttt{<p>} com o preço;
		\item confira três \texttt{<article>} no \texttt{Ctrl+U};
		\item acrescente um quarto produto ao array e confira que ele aparece.
	\end{enumerate}

	\vspace{0.2cm}
	Tempo: 15 minutos.
\end{frame}
}

{\estiloresposta
\begin{frame}[fragile]{Resposta: o catálogo com \texttt{foreach}}
	\begin{minted}[fontsize=\scriptsize]{html+php}
<body>
  <p><?= count($produtos) ?> produtos</p>
  <?php foreach ($produtos as $produto): ?>
    <article class="cartao">
      <h2><?= $produto["nome"] ?></h2>
      <p><?= $produto["descricao"] ?></p>
      <p>R$ <?= number_format($produto["preco"], 2, ",", ".") ?></p>
    </article>
  <?php endforeach; ?>
</body>
	\end{minted}

	\vspace{0.2cm}
	O quarto produto aparece sem mudança no HTML: o \texttt{foreach} gera um
	cartão por elemento do array.
\end{frame}
}

\begin{frame}[fragile]{Funções}
	\begin{minted}{php}
function formataPreco(float $valor): string {
  return "R$ " . number_format($valor, 2, ",", ".");
}

echo formataPreco(27.625);   // R$ 27,63
	\end{minted}

	\begin{itemize}
		\item tipos do parâmetro (\texttt{float \$valor}) e do retorno
		      (\texttt{: string}) são opcionais;
		\item com eles, \texttt{formataPreco("abc")} interrompe o script com
		      \texttt{TypeError};
		\item variáveis de fora não são visíveis dentro da função: o que ela usa
		      entra pelos parâmetros.
	\end{itemize}
\end{frame}

{\estilopratica
\begin{frame}[fragile]{Prática: o desconto com \texttt{if}}
	No \texttt{catalogo.php}:

	\begin{enumerate}
		\item escreva a função \texttt{formataPreco} no bloco do topo, e use-a em
		      todos os preços;
		\item troque o parágrafo do preço por um \texttt{if} na sintaxe
		      alternativa: com desconto, o preço original dentro de \texttt{<s>} e
		      o preço com desconto ao lado; sem desconto, só o preço.
	\end{enumerate}

	\vspace{0.2cm}
	Confira:

	\begin{minted}{text}
Café em grão     R$ 32,50 R$ 27,63
Chá mate         R$ 12,90
Pão de mel       R$ 8,00 R$ 7,20
	\end{minted}

	Tempo: 15 minutos.
\end{frame}
}

{\estiloresposta
\begin{frame}[fragile]{Resposta: o desconto com \texttt{if}}
	\begin{minted}[fontsize=\scriptsize]{html+php}
<?php
function formataPreco(float $valor): string {
  return "R$ " . number_format($valor, 2, ",", ".");
}
// ... array $produtos ...
?>
<!-- dentro do foreach, no lugar do parágrafo do preço -->
<?php if ($produto["desconto"] > 0): ?>
  <p>
    <s><?= formataPreco($produto["preco"]) ?></s>
    <?= formataPreco($produto["preco"] * (1 - $produto["desconto"])) ?>
  </p>
<?php else: ?>
  <p><?= formataPreco($produto["preco"]) ?></p>
<?php endif; ?>
	\end{minted}
\end{frame}
}

% ==========================================
\section{JSON}
% ==========================================

\begin{frame}[fragile]{\texttt{json\_decode}}
	O mesmo \texttt{produtos.json} lido pelo \texttt{fetch}, agora no servidor:

	\begin{minted}{php}
$conteudo = file_get_contents("produtos.json"); // texto
$produtos = json_decode($conteudo, true);        // array
	\end{minted}

	\vspace{0.2cm}
	\begin{itemize}
		\item \texttt{true}: objetos do JSON viram arrays associativos;
		\item o resultado tem a forma do array de arrays, e o mesmo
		      \texttt{foreach} serve;
		\item \texttt{json\_encode} faz o caminho inverso, e é como um back-end
		      PHP responde a um \texttt{fetch}.
	\end{itemize}
\end{frame}

\begin{frame}[fragile]{JSON com erro}
	Com erro de sintaxe no arquivo, \texttt{json\_decode} devolve \texttt{null},
	sem mensagem:

	\begin{minted}{php}
$produtos = json_decode('[{"id": 1,}]', true);
var_dump($produtos);           // NULL
echo json_last_error_msg();    // Syntax error
	\end{minted}

	\vspace{0.2cm}
	O \texttt{foreach} logo depois avisa:

	\begin{minted}[fontsize=\scriptsize]{text}
Warning: foreach() argument must be of type array|object, null given
	\end{minted}
\end{frame}

{\estilopratica
\begin{frame}[fragile]{Prática: o catálogo pelo \texttt{produtos.json}}
	\begin{enumerate}
		\item copie para \texttt{aula-php/} o \texttt{produtos.json} da Entrega 2,
		      ou escreva os três produtos em JSON;
		\item abra \texttt{http://localhost:8000/produtos.json} e confira que o
		      servidor entrega o arquivo;
		\item no \texttt{catalogo.php}, troque o array escrito no código pela
		      leitura do arquivo;
		\item acrescente um produto ao JSON e confira que ele aparece sem
		      alteração no PHP.
	\end{enumerate}

	\vspace{0.2cm}
	\texttt{Undefined array key "desconto"}: falta essa chave em algum produto do
	arquivo. Acrescente \texttt{"desconto": 0}, ou ajuste o nome no PHP.

	\vspace{0.2cm}
	Tempo: 10 minutos.
\end{frame}
}

{\estiloresposta
\begin{frame}[fragile]{Resposta: o catálogo pelo \texttt{produtos.json}}
	\begin{minted}{php}
<?php
function formataPreco(float $valor): string {
  return "R$ " . number_format($valor, 2, ",", ".");
}

$conteudo = file_get_contents("produtos.json");
$produtos = json_decode($conteudo, true);
?>
	\end{minted}

	\vspace{0.2cm}
	O HTML abaixo do bloco não muda.

	\vspace{0.2cm}
	Na Entrega 3, o lugar do \texttt{json\_decode} é ocupado pela consulta ao
	banco de dados.
\end{frame}
}

{\estilopratica
\begin{frame}[fragile]{Prática: filtro por categoria}
	No \texttt{catalogo.php}:

	\begin{enumerate}
		\item crie no topo a variável \texttt{\$categoria}, com o valor
		      \texttt{"Bebidas"};
		\item dentro do \texttt{foreach}, mostre apenas os produtos dessa
		      categoria, com um \texttt{if};
		\item no topo, conte em \texttt{\$mostrados} os produtos da categoria,
		      com outro \texttt{foreach} e um \texttt{if}, e use a variável no
		      \texttt{<p>} da quantidade;
		\item troque o valor de \texttt{\$categoria} e recarregue.
	\end{enumerate}

	\vspace{0.2cm}
	Tempo: 10 minutos.
\end{frame}
}

{\estiloresposta
\begin{frame}[fragile]{Resposta: filtro por categoria}
	\begin{minted}[fontsize=\scriptsize]{html+php}
<?php
// ... formataPreco e json_decode ...
$categoria = "Bebidas";
$mostrados = 0;
foreach ($produtos as $produto) {
  if ($produto["categoria"] === $categoria) {
    $mostrados++;
  }
}
?>
<!-- ... -->
<p><?= $mostrados ?> produtos</p>
<?php foreach ($produtos as $produto): ?>
  <?php if ($produto["categoria"] === $categoria): ?>
    <article class="cartao"> ... </article>
  <?php endif; ?>
<?php endforeach; ?>
	\end{minted}

	\small
	Na aula de formulários, \texttt{\$categoria} passa a vir do endereço,
	como em \texttt{catalogo.php?categoria=Bebidas}.
\end{frame}
}

\begin{frame}{Quando dá errado}
	\scriptsize
	\begin{tabular}{>{\raggedright\arraybackslash}p{0.48\textwidth}>{\raggedright\arraybackslash}p{0.42\textwidth}}
		\toprule
		Sintoma & Causa provável \\
		\midrule
		página mostra o código PHP, ou o navegador baixa o arquivo & sem o
		servidor do PHP \\ \midrule
		página em branco & \texttt{display\_errors} desligado; erro no
		terminal \\ \midrule
		\texttt{Parse error: syntax error, unexpected token "echo"} &
		\texttt{;} faltando na linha anterior \\ \midrule
		\texttt{Undefined variable} & nome digitado diferente, ou variável de
		fora usada na função \\ \midrule
		\texttt{Undefined array key} & chave inexistente no elemento \\
		\midrule
		\texttt{Unsupported operand types: string + float} & \texttt{+} no
		lugar de \texttt{.} \\ \midrule
		\texttt{Array} na página & \texttt{echo} de um array inteiro \\
		\midrule
		\texttt{\$nome} literal na página & aspas simples \\ \midrule
		\texttt{foreach() argument must be of type array|object, null given} &
		JSON com erro ou caminho errado \\
		\bottomrule
	\end{tabular}
\end{frame}

\begin{frame}{Exercícios e próximos passos}
	Os exercícios desta aula não têm entrega e não valem nota. O enunciado
	completo está na Parte prática do material de PHP.

	\vspace{0.3cm}
	Próxima aula: formulários. \texttt{\$\_GET} e \texttt{\$\_POST}, validação no
	servidor, \texttt{htmlspecialchars} e \texttt{include}.

	\vspace{0.3cm}
	A Entrega 3 do trabalho (PHP e MySQL) vence em 27/11, às 23h59. Não há
	entrega em atraso: vale o último \emph{commit} enviado até o prazo.
\end{frame}

\end{document}
